\documentclass[aspectratio=169]{beamer}

% Packages
\usepackage[utf8]{inputenc}
\usepackage[english]{babel}

% Theme and colors
\usetheme{Madrid}
\usecolortheme{default}

% Custom colors
\definecolor{chatgptgreen}{RGB}{16, 163, 127}
\setbeamercolor{structure}{fg=chatgptgreen}
\setbeamercolor{title}{fg=white, bg=chatgptgreen}
\setbeamercolor{frametitle}{fg=white, bg=chatgptgreen}

% Title page information
\title{Meet ChatGPT}
\subtitle{Your New Digital Assistant}
\author{}
\date{}

\begin{document}

\begin{frame}
    \titlepage
\end{frame}

\begin{frame}{Slide 1: What is ChatGPT?}
    \begin{itemize}
        \item \textbf{Concept:} Imagine having a very smart, well-read assistant who is awake 24/7 and ready to help you with almost anything. It is a computer program that you can talk to just like a human.
        \item \textbf{Real World Example:} Instead of searching Google and reading 10 different websites to find an answer, you just ask ChatGPT a question, and it writes the exact answer for you immediately.
        \item \textbf{Prompt to try:} \textit{"I am planning a 3-day trip to Vienna. I love classical music and visiting old cafes. Can you write a day-by-day itinerary for me?"}
    \end{itemize}
\end{frame}

\begin{frame}{Slide 2: Keeping Things Fresh (Why New Chats Matter)}
    \begin{itemize}
        \item \textbf{Concept:} Just like humans, if you talk about too many different things in one long conversation, ChatGPT can get "exhausted" or confused. It might mix up your previous topics.
        \item \textbf{The Rule:} Always start a \textbf{New Chat} for a new topic. If you were talking about a recipe, and now want to talk about a rock band, hit the "New Chat" button!
        \item \textbf{Real World Example:} If you have a chat where you are asking for help writing a formal email to your boss, don't suddenly ask it for a pizza recipe in that same chat. Start a new one!
    \end{itemize}
\end{frame}

\begin{frame}{Slide 3: Your Personal Translator}
    \begin{itemize}
        \item \textbf{Concept:} You can turn ChatGPT into a dedicated, lightning-fast translator.
        \item \textbf{How to do it:} Start a new chat and give it one single instruction at the beginning.
        \item \textbf{Prompt to copy and paste:} \textit{"From now on, whatever I type in Ukrainian, translate it to English. And whatever I type in English, translate it to Ukrainian. Do not add any other text, just give me the translation."}
        \item \textbf{The Magic:} You never have to ask it to translate again in that specific chat. Just open that chat, type your text, and it will automatically translate it every time!
    \end{itemize}
\end{frame}

\begin{frame}{Slide 4: Searching the Web \& Finding Links}
    \begin{itemize}
        \item \textbf{Concept:} ChatGPT isn't just a brain; it can also search the live internet for you.
        \item \textbf{Real World Example:} You want to buy tickets for a concert, but you don't know where to look or what the official website is.
        \item \textbf{Prompt to try:} \textit{"Find me the official website to buy tickets for the upcoming Coldplay concert in Warsaw, and give me the direct link."}
    \end{itemize}
\end{frame}

\begin{frame}{Slide 5: Seeing the World Through Pictures}
    \begin{itemize}
        \item \textbf{Concept:} ChatGPT has "eyes"! You can upload photos from your phone and ask questions about them.
        \item \textbf{Real World Example:} You are walking in the park and see a beautiful, strange flower, or you are in a store and see a gadget you don't understand.
        \item \textbf{How to do it:} Snap a picture, upload it to ChatGPT, and ask:
        \item \textbf{Prompt to try:} \textit{"What kind of flower is this in the picture? How often does it need to be watered if I keep it in my apartment?"}
    \end{itemize}
\end{frame}

\begin{frame}{Slide 6: A Listening Ear (Psychological Support)}
    \begin{itemize}
        \item \textbf{Concept:} Sometimes we all feel a little sad, stressed, or just need someone to talk to.
        \item \textbf{Real World Example:} You had a very stressful day at work and you just want to vent, or you are feeling anxious about an upcoming event.
        \item \textbf{Prompt to try:} \textit{"I had a really terrible day today and I feel very stressed. Can you just listen to me and offer some comforting words?"}
    \end{itemize}
\end{frame}

\begin{frame}{Slide 7: Everyday Helper (Cooking, Health, and a Warning!)}
    \begin{itemize}
        \item \textbf{Concept:} It’s a fantastic assistant for daily life. You can tell it what ingredients are in your fridge, and it will give you a step-by-step recipe.
        \item \textbf{Real World Example:} You open your fridge and only see chicken, rice, and some carrots.
        \item \textbf{Prompt to try:} \textit{"I only have chicken breast, white rice, and carrots in my kitchen. What is a simple and tasty dinner I can make with just these ingredients?"}
        \item \textbf{The Warning:} While it can give general health or wellness advice, \textbf{never trust AI blindly with serious matters}. It is not a real doctor. Always verify important health, medical, or safety information with a human professional.
    \end{itemize}
\end{frame}

\begin{frame}{Slide 8: The Ultimate Teacher (Math, Science \& Tech)}
    \begin{itemize}
        \item \textbf{Concept:} ChatGPT is brilliant at explaining complicated things in a very simple way.
        \item \textbf{Real World Example:} You hear people talking about "Blockchain" or "Quantum Physics" on the news, and you have no idea what it means.
        \item \textbf{Prompt to try:} \textit{"Explain to me what a Black Hole is, but explain it using simple words like I am a 10-year-old child."}
    \end{itemize}
\end{frame}

\begin{frame}{Slide 9: Let's Try It!}
    \begin{itemize}
        \item \textbf{Concept:} The best way to learn is by doing. Let's open the app and try one of the prompts right now!
    \end{itemize}
\end{frame}

\end{document}